
\subsubsection{Multi-Objective Code Generation}

Recent work in code generation extends beyond single-objective correctness optimization. PPOCoder (Shojaee et al., 2023) pioneered execution-based reinforcement learning using PPO with compilation feedback. CodeRL (Le et al., 2022) similarly leverages RL with execution feedback for program synthesis.

Production systems require optimizing multiple objectives. The standard approach uses weighted reward combinations with manual tuning (Srivastava \& Aggarwal, 2025; ReTool, Feng et al., 2025). Practitioners tune weight vectors through grid search or Bayesian optimization.

Recent work proposes architectural factorization as an alternative. Wong \& Tan (2025) apply RLHF with crowd-sourced feedback using Bayesian optimization to integrate preferences. PairCoder (Zhang et al., 2024) introduces navigator-driver agent collaboration with feedback-driven refinement. These methods implicitly assume quality objectives can be decomposed into architectural modules—an assumption not validated empirically.

\subsubsection{Multi-Task Learning}

Architectural factorization for code draws from multi-task learning where task-specific subspaces with orthogonality constraints prove effective (Caruana, 1997; Ruder, 2017). Pareto Multi-Task Learning (Lin et al., 2019) and PCGrad (Yu et al., 2020) address task conflicts through gradient balancing, assuming separable task structures exist.

The disentanglement literature (Bengio et al., 2013; β-VAE) seeks representations where latent dimensions correspond to interpretable factors. Success stories come from domains with known ground-truth factors (image rotation, color, scale). Whether code quality aspects have analogous structure remains an empirical question.

\subsubsection{Gap in Validation}

Prior multi-objective code generation work proceeds directly from architectural intuition to system design without validating empirical separability. Our contribution provides the validation methodology this literature lacks.
